\documentclass[11pt,fontset=fandol]{ctexart}
\usepackage[a4paper,margin=23mm]{geometry}
\usepackage{amsmath,amssymb,mathtools,bm,mathrsfs,esint}
\usepackage{graphicx,longtable,booktabs,array,calc}
\usepackage[normalem]{ulem}
\usepackage{xcolor,hyperref}
\hypersetup{hidelinks}
\setcounter{secnumdepth}{-1}
\newcounter{none}
\providecommand{\tightlist}{\setlength{\itemsep}{0pt}\setlength{\parskip}{0pt}}
\providecommand{\passthrough}[1]{#1}
\usepackage{amssymb,mathrsfs,esint,bm,newunicodechar}
\allowdisplaybreaks
\setlength{\emergencystretch}{3em}
\providecommand{\llangle}{\langle\!\langle}
\providecommand{\rrangle}{\rangle\!\rangle}
\newunicodechar{∞}{\ensuremath{\infty}}
\newunicodechar{∈}{\ensuremath{\in}}
\newunicodechar{∑}{\ensuremath{\sum}}
\newunicodechar{∫}{\ensuremath{\int}}
\newunicodechar{≤}{\ensuremath{\leq}}
\newunicodechar{≥}{\ensuremath{\geq}}
\newunicodechar{≠}{\ensuremath{\neq}}
\newunicodechar{→}{\ensuremath{\to}}
\newunicodechar{⇒}{\ensuremath{\Rightarrow}}
\newunicodechar{α}{\ensuremath{\alpha}}
\newunicodechar{β}{\ensuremath{\beta}}
\newunicodechar{γ}{\ensuremath{\gamma}}
\newunicodechar{λ}{\ensuremath{\lambda}}
\newunicodechar{φ}{\ensuremath{\varphi}}
\newunicodechar{ε}{\ensuremath{\varepsilon}}
\newunicodechar{σ}{\ensuremath{\sigma}}
\newunicodechar{θ}{\ensuremath{\theta}}
\newunicodechar{δ}{\ensuremath{\delta}}
\newunicodechar{η}{\ensuremath{\eta}}
\newunicodechar{π}{\ensuremath{\pi}}
\newunicodechar{ξ}{\ensuremath{\xi}}
\newunicodechar{Σ}{\ensuremath{\Sigma}}
\newunicodechar{Δ}{\ensuremath{\Delta}}
\newunicodechar{∂}{\ensuremath{\partial}}
\newunicodechar{∀}{\ensuremath{\forall}}
\newunicodechar{∃}{\ensuremath{\exists}}
\newunicodechar{√}{\ensuremath{\surd}}
\newunicodechar{−}{\ensuremath{-}}
\newunicodechar{μ}{\ensuremath{\mu}}
\newunicodechar{ρ}{\ensuremath{\rho}}
\newunicodechar{τ}{\ensuremath{\tau}}
\newunicodechar{ω}{\ensuremath{\omega}}
\newunicodechar{Ω}{\ensuremath{\Omega}}
\newunicodechar{χ}{\ensuremath{\chi}}
\newunicodechar{ψ}{\ensuremath{\psi}}
\newunicodechar{Γ}{\ensuremath{\Gamma}}
\newunicodechar{Φ}{\ensuremath{\Phi}}
\newunicodechar{Λ}{\ensuremath{\Lambda}}
\newunicodechar{ℝ}{\ensuremath{\mathbb{R}}}
\newunicodechar{ℤ}{\ensuremath{\mathbb{Z}}}
\newunicodechar{ℕ}{\ensuremath{\mathbb{N}}}
\newunicodechar{△}{\ensuremath{\triangle}}
\newunicodechar{⊗}{\ensuremath{\otimes}}
\begin{document}
\section{2023年杭州电子科技大学601数学分析考研真题}\label{ux5e74ux676dux5ddeux7535ux5b50ux79d1ux6280ux5927ux5b66601ux6570ux5b66ux5206ux6790ux8003ux7814ux771fux9898}

来源：\texttt{HDU/601数学分析/2023/2023年杭州电子科技大学601数学分析考研真题.pdf}。原
PDF 页序：1, 2。

\section{杭州电子科技大学2023年数学分析真题}\label{ux676dux5ddeux7535ux5b50ux79d1ux6280ux5927ux5b662023ux5e74ux6570ux5b66ux5206ux6790ux771fux9898}

\subsection{一.填空题 (本大题共8 小题,每小题 5 分，共 40
分)}\label{ux4e00.ux586bux7a7aux9898-ux672cux5927ux9898ux51718-ux5c0fux9898ux6bcfux5c0fux9898-5-ux5206ux5171-40-ux5206}

1.已知 \(f ( x ) = { \frac { x } { \sqrt { 1 + x ^ { 2 } } } }\) ，设
\(f _ { n } ( x ) = f ( f ( \cdots f ( x ) ) )\) ，则f2023（0)=

\begin{enumerate}
\def\labelenumi{\arabic{enumi}.}
\setcounter{enumi}{1}
\tightlist
\item
  \(\lim _ { n \to \infty } n ( { \sqrt { n ^ { 2 } + 1 } } - n ) =\)
\end{enumerate}

\(F ( x ) = \int _ { a } ^ { x } f ( t ) ( t - x ) d t\)
则F(x)在{[}a,b{]}上的 导数F(x)= 考研考傅 贩必究

\begin{enumerate}
\def\labelenumi{\arabic{enumi}.}
\setcounter{enumi}{3}
\tightlist
\item
  \(\int _ { - \infty } ^ { + \infty } { \tt e } ^ { - x ^ { 2 } } d x =\)
\end{enumerate}

5.函数 W
\(f ( x ) = \left\{ \begin{array} { l l } { \displaystyle \cos \frac { \pi x } { 2 } , \big | x \big | \leq 1 } \\ { \big | x - 1 \big | , \big | x \big | > 1 } \end{array} \right.\)
的间断点为

6.级数
\(\sum _ { n = 2 } ^ { \infty } ( 1 + { \frac { 1 } { n } } ) ^ { n ( n + 1 ) } x ^ { n }\)
的收敛半径为考

7.设曲线L是圆 \(x ^ { 2 } + y ^ { 2 } = r ^ { 2 }\)
在第一象限内的部分，则 \(\int _ { L } x y d l =\) 研

\begin{enumerate}
\def\labelenumi{\arabic{enumi}.}
\setcounter{enumi}{7}
\tightlist
\item
  \(\int { \frac { 1 } { x ( 1 + x ^ { n } ) } } d x =\) ·
\end{enumerate}

\subsection{二.计算题(本大题共5小题，每小题10分，共50
分)}\label{ux4e8c.ux8ba1ux7b97ux9898ux672cux5927ux9898ux51715ux5c0fux9898ux6bcfux5c0fux989810ux5206ux517150-ux5206}

1.计算
\(\int _ { L } { ( - 2 x e ^ { - x ^ { 2 } } \sin y ) d x } + { ( e ^ { - x ^ { 2 } } \cos y + x ^ { 6 } ) d y }\)
,这里L为从点(1,0)到(-1,0) 的半圆
\(y = \sqrt { 1 - x ^ { 2 } } ( - 1 \leq x \leq 1 )\) 無考博 必究

2.求
\(\lim _ { x \to \infty } \left[ x - x ^ { 2 } \ln \left( 1 + { \frac { 1 } { x } } \right) \right]\)

3.求 \(\sum _ { n = 1 } ^ { \infty } { \frac { 2 n } { 3 ^ { n } } }\)
的和.

4.计算 \(\oint _ { \Gamma } x ^ { 2 } y d x + y ^ { 2 } d y + z d z\)
,其中r是圆柱 \(x ^ { 2 } + y ^ { 2 } = 1\) 与平面x+z=1的交线，从z
轴的正向看去，为逆时针方向.

5.计算
\(\iint _ { \Omega } { x ^ { 2 } \sqrt { { x ^ { 2 } + y ^ { 2 } } } d x d y d z }\)
,其中Ω是曲面 \(z = \sqrt { x ^ { 2 } + y ^ { 2 } }\) 与平面
\(z = x ^ { 2 } + y ^ { 2 }\) 所围成有 学考研考 版必究的区域. 老博

\subsection{三.证明题(本大题共4小题，每小题10分，共40
分)}\label{ux4e09.ux8bc1ux660eux9898ux672cux5927ux9898ux51714ux5c0fux9898ux6bcfux5c0fux989810ux5206ux517140-ux5206}

1.证明：
\(\sum _ { n = 1 } ^ { \infty } ( - 1 ) ^ { n } { \frac { x ^ { 2 } + n } { n ^ { 2 } } }\)
在任何有限区域{[}a,b{]}上一致收敛，而在实轴上任何一点都不绝对收敛.
版权所

2.设奇函数f(x)在{[}-1,1{]}上具有二阶导数，且f(1)=1,证明：

(1)存在 \(\xi \in ( 0 , 1 )\) ，使得f(§)=1;

(2)存在η∈(-1,1)，使得
\(f ^ { \prime \prime } ( \eta ) + f ^ { \prime } ( \eta ) = 1\)

3.设f(x)在{[}0,1{]}上连续，证明：
\(\iint f ( 1 - y ) f ( x ) d x d y = \frac { 1 } { 2 } \Bigl [ \int _ { 0 } ^ { 1 } f ( x ) d x \Bigr ] ^ { 2 }\)
其中D

为以O(0,0),A(0,1),B(1,0)为顶点的三角形所围成的区域.

4.已知△表示L2pl2ce算符，即：
\(\Delta u = \frac { \hat { \sigma } ^ { 2 } u } { \hat { \sigma } { x } ^ { 2 } } + \frac { \hat { \sigma } ^ { 2 } u } { \hat { \sigma } { y } ^ { 2 } } + \frac { \hat { \sigma } ^ { 2 } u } { \hat { \sigma } { z } ^ { 2 } }\)
．证明：

\[
\begin{array} { r } { \Delta ( u \nu ) = ( \Delta u ) \nu + 2 \nabla u \cdot \nabla \nu + u ( \Delta \nu ) . } \end{array}
\]

这里Vu表示函数u在点(x,y,z)处的梯度，``.''表示点乘.

\subsection{四.综合题(本大题共2小题，共30分)}\label{ux56db.ux7efcux5408ux9898ux672cux5927ux9898ux51712ux5c0fux9898ux517130ux5206}

1.(10分)求两个曲面x+2y=1和 \(x ^ { 2 } + 2 y ^ { 2 } + z ^ { 2 } = 1\)
的交线上距离原点最近的点.

2.设
\(f ( x ) = \int _ { 1 } ^ { + \infty } \frac { \cos ( x t ) } { t ( 1 + t ^ { 2 } ) } d t , x \in ( - \infty , + \infty )\)

(1）试证明积分
\(\int _ { 1 } ^ { + \infty } \frac { \cos ( x t ) } { t ( 1 + t ^ { 2 } ) } d t\)
关于x在 \(( - \infty , + \infty )\) 上一致收敛.

(2）试证明函数f(x)在 \(( - \infty , + \infty )\) 上一致连续.
\end{document}
